\chapter*{Prologue: The Night of the Kettle}
\addcontentsline{toc}{chapter}{Prologue: The Night of the Kettle}
\markboth{PROLOGUE}{PROLOGUE}

Dusk over Hollowmere, and the lamps going up one by one, like the city buttoning its coat.

From the ridge of the Fish-Giver's roof I can see most of what matters: the chimney forests of the Shambles, the canals, the fog beginning its evening climb up from the Mere the way it always does. Somewhere below, a door slams. Somewhere else, a man says something unkind about a horse. The city, in other words, in working order.

She puts the kettle on when the lamps go up. I am punctual. She is predictable. The arrangement works, and has worked for fifteen years, which is a long time in a cat.

Fifteen. It means I am stiff in the mornings until the sun gets into me, it means I take the stairs and not the banister, and it means I retired some while ago without telling anybody, there being nobody to tell. My working nights are behind me. My employment now is to be somewhere warm at the correct time, and I am extremely good at it.

Her name, for the record --- because tonight of all nights the record matters --- is Elspeth Vane. The city calls her the Needle of Hollowmere, mostly behind her back, always with respect. She is a witch, which in this city is a trade like any other: she sews. Wards, hems, bindings, the long invisible seams that keep a place from coming apart at its corners. Magic here is needlework. Anyone who tells you otherwise is selling candles.

And I am Ash. Her cat. Other titles have been proposed over the years: familiar, companion, mouser-in-residence, \emph{that animal}. I answer to none of them and come to all of them, provided there is a fire on.

Tonight calls for a gift. A vole, ideally. Voles say what flowers cannot.

\scenebreak

There. Under the chimney pot. Fat, smug, unaware.

I go low. The tiles are still day-warm under my belly, and the wind is doing me the courtesy of blowing the wrong way for the vole and the right way for me. There is a red feeling that lives at the bottom of a cat, older than cities, and it comes up through me now, the world narrowing to one point with fur on it. It is the one part of me that has not aged.

The rest of me has. I wait longer than I once needed to, because waiting is cheap and the other thing is not. I let it come two tiles nearer than a young cat would have troubled to.

Then the vole makes its argument, which is brief, and mostly consists of running. I make mine. Mine is better, by about a whisker. There is a scuffle among the chimney pots which I will describe, in the interest of dignity, as \emph{professional}, and then it is over, and the vole has been acquired. Thoughtfully dead. Arranged neatly, carried gently.

I sat afterward longer than the work deserved, and did not mention it to anybody.

She will say ``not on the table,'' and mean \emph{thank you}. She has said it every time. I have put it on the table every time.

I take the rooftop road toward home with the vole warm in my mouth and the evening going purple at its edges. There is a gap over the alley at Cooper's that I used to cross without consulting anybody, including myself. I stood at the edge of it and looked, and counted the four seconds she taught me, and went round by the wall instead.

I spent that hour thinking about a sunbeam I remember from March. I want that noted. Whatever else the record says of me, let it say I was happy.

\scenebreak

Below me now: the Shambles, shut for the night, stalls sleeping under canvas. Eight bells sound from the tower, the whole count, one to eight, dropped into the fog like coins into a well.

There used to be a ninth bell. There is no ninth any more. Nobody remembers when that changed, and nobody finds it strange, and I have never met anyone but me who noticed there was a gap where a sound should be.

And between two chimneys, strung across the dark like a hair caught in a comb: a thread of silver light. Humming.

Everything alive trails a thread, if you have the eyes for it. They run through this city out of windows, along gutters, wound around wrists and doorposts and gateposts. They are the underlay of everything, which she calls the Weft. Witches stitch with them. Cats only see them. Hers I know by heart: silver, with a particular twist to it, like handwriting.

This one is silver too. And it is nobody's I know.

I did not chase the silver thread. A cat of my age, carrying a gift, does not go haring off across the roofline after every unexplained glimmer. I looked at it. I considered it. I made a note in the ledger of things that are wrong tonight and went on my way, and the lie I told myself about not \emph{wanting} to chase it weighed nothing at all.

Almost nothing.

\scenebreak

Street level, Needle Lane, and the fog rolls up from the Mere to meet me, and the shadows deepen, deep enough to wear, if you know how. I have been a rumor most of my life.

Something small and hungry climbs out of the gutter.

A gutter-wisp: a swallowed lantern of a thing, all glow and appetite. They eat small lies, and the city breathes out enough of those to keep a great many of them fed. This one heard mine. The one about not wanting to chase the thread. It has come to collect.

I would like to tell you I handled it with grace. What happened is that a cat and a lit appetite had a disagreement in a foggy gutter, and the fog was on my side, and the wisp went back down the drain, dimmer.

I sat longer than I meant to afterward, putting my coat back in order, and that was when I noticed the quiet.

No pigeons.

Understand: pigeons are never quiet. The roosts over Needle Lane keep up a continuous session of complaint and scandal from dusk until dawn, and I have slept under that noise my whole life. It was not there. The roofline sat against the sky, empty as a shelf.

\scenebreak

The lamps on Needle Lane are lit in the wrong order.

You would have to live here to see it. The Lamplighters walk a sequence, every night: third lamp, then first, then the one by the pump. The order is the point. I have heard one of them say so, on a doorstep, to an apprentice: \emph{we light them in an order, lad, and the order is the stitch}. Tonight the lane is lit dark-bright-dark-dark-bright, a signature forged by someone who has only ever seen it written.

The Lamplighters never miss. The dog at Number Twelve is awake, I can smell him being awake, and he does not bark, which he has never once in his loud and tedious life failed to do.

And halfway down the lane, the fog stops being fog.

It has sleeves. It has a collar. It has the shape of a person who is not in it. And that is when I see what has been done to the lane: window-threads, door-threads, the little tethers of household luck, all of them \emph{cut}, and sliding through the dark toward this one standing shape like water finding a drain, and the shape gathering them in.

They have a name, things like this. Rag-wraiths: clothes whose owners are missing. Not dead --- \emph{missing}. The distinction matters to the things themselves, I am told.

It turns what it has instead of a face toward me. Whatever unpicked these lamps is still at its work, and I am standing in its light.

I fought it. I will say this once, plainly: I went in with everything I had left in me, and then with the thing of hers I carry, the pale cold weight under my ribs that is not mine and never was. I felt it go out of me and do nothing. That is a feeling I had no name for then.

It was above my weight and it knew it. Ten years ago it would still have been above my weight. In the end I did the intelligent thing, which the stories never put in: I slipped away. Three streets before I stopped, and I stopped because my legs decided it, not me, and I stood in a doorway with my sides going like a bellows and waited to be a cat again.

Nothing followed. Nothing was missing.

And yet. It was not hunting me. All that gathered-up stolen thread, and it spent none of it on a cat. It was \emph{working}, and I did not know yet at what.

\scenebreak

Her window is ahead now.

The shutters are closed.

I stood in the road and could not make that mean anything. They are old green shutters and they fold back against the brick, and they have been folded back every evening of my life: in rain, in frost, in the week the chimney caught. She did not close them for weather. She did not close them for company. She did not close them.

Something in me tried to file it and could not find the drawer. So I did what anybody does with a fact that will not sit down. I told myself she was out.

Two ways home from here. The long way passes the Lamplighters' hall, and whatever silenced the lamps went wrong there first. The fast way is over the fences.

I took the long way. Eleven minutes, perhaps twelve. Because the route had a wrongness at the head of it, and because her voice said \emph{look before you leap, Ash, it is the entire difference between brave and stupid}, and because the fences are a young cat's road and I had already had my breath taken off me once that night. All three are true. I have had a long time to decide which one I was obeying, and I have not decided.

The hall stood open. Door wide, ladders racked, coats on their pegs. The kettle on their little stove was cold, and a guild hall's kettle is never cold. Forty men and women walk out of this room every evening and light a city. Nobody was in it.

Caught on a lamp-hook by the door, glinting where it had snagged: a thread of silver.

\emph{Her} stitching. I would know it under floodwater. She was here --- tonight, in this room, my witch, who has not walked out on an evening in years. Why was she \emph{here}?

I ran then. Out of the hall and eleven streets home, and I do not remember eight of them. The fog thinning near the Mere. Ordinary houses, ordinary gates, a cat crossing an ordinary city at a speed that was going to cost him. Her street. Her gate. Her drainpipe, the one I have used nine hundred times, and up --- slower than nine hundred times, my back legs arguing the whole way --- and in through the workroom window ---

and the house is doing the wrong thing with sound.

A kettle, screaming. And nobody going to it.

\scenebreak

The workroom first, and the workroom is wrong. The threads on the walls, forty years of her work hung like a garden, cut. Every one of them snipped at the same height, cleanly, as if by someone \emph{tidying}.

The parlor door stands open, and I do not go through it. Not yet. I stop at the threshold, low, and I look.

Something is standing in the parlor doorway on the far side, wearing the room.

It is person-shaped the way a coat on a hook is person-shaped, and it is made of \emph{her} threads: it has been pulling forty years of her work out of the walls and putting it \emph{on}. Where the light of the hearth should catch a face there is a weave, still moving, still settling, thread finding thread with a sound like the smallest rain.

That is the part I will keep, of everything that night: it looked at me, whatever it looks with, and it \emph{decided I was not an interruption}, and it went back to its work.

Every hair I have is standing up. Good. Let them.

I went for it. Of course I went for it. She was somewhere past that doorway and the thing between us was wearing her wards. I remember the first spring off the floorboards, and that it was a poor one, and that I made it anyway. I remember the weave turning, unhurried, and a sleeve coming around like a laden clothesline let go in a gale. I remember the cold weight under my ribs going out of me a second time that night, all of it at once, everything of hers I had left to spend, and it was not enough, and the room went sideways, and then the cold came up through the floor in a way that floors do not do.

I want to tell you dying hurts. It would make better telling. What it does is \emph{leave}: the cold goes first, then the noise, then the part of you that was worried about the noise. What is left over found itself standing in a corridor that was all filing and no building, in a light the color of four in the morning.

Nothing hurt. I found this deeply suspicious.

A desk. A lamp. A queue of one, and the one was me.

The being behind the desk was thin the way a ruled line is thin, translucent at the cuffs, with spectacles that had given up asking to be needed. He did not look up. Around him, shelves went away in every direction that shelves can go, and some that they cannot, all of them full, every folder tied with grey ribbon.

``Name?''

``Ashcat,'' I said. ``Of Needle Lane.'' It seemed a night for the full formal.

``Needle Lane.'' Something moved behind the spectacles that was almost interest. ``Yes. We have had a busy evening from Needle Lane.''

He wrote. The pen made the sound of winter. Then he reached the column that mattered, and stopped.

``Nine allotted.'' He said it to the paper, not to me. He turned the sheet over, as though the back of it might explain; the back of it did not. ``I have not filed a nine since I was at a much smaller desk. Vane's familiar.'' A pause you could have hung a coat on. ``One spent.''

He did not say whose the busy evening was. I did not ask. Every whisker I own said this was a room where questions were expensive.

``The arrangement, since it is your first,'' said the Clerk. ``Deaths are appointments. Tonight, you had one.'' He said it without cruelty, the way a man reads weather. ``Arrive dead \emph{without} an appointment and you will be refused at this desk --- we are a court, not a revolving door. A true death --- the kind you will always see coming --- costs a life, and after the first, the Toll. Any other defeat is refused and returned. The first visit is complimentary. Do consider it a courtesy rather than a habit.''

The stamp came down. It was, I noticed, a paw print. Neither of us remarked on this.

The stairs up were longer than the fall down.

\scenebreak

Hearing came back first. Floorboards under one cheek. The smell of cold hearth. Four paws --- count them --- four. Whatever had been standing in the doorway was not standing in the doorway.

Get up, Ash.

The parlor.

Elspeth in her chair. The light gone out of her hands.

\begin{center}\ldots\end{center}

The kettle was still screaming. It had been screaming this whole time, and I did not hear it until now. She hates it when it does that.

I took it off the fire. One paw, on the handle, the way I had watched her do it ten thousand times from the hearthrug. It took three tries. Nobody will ever know.

And then the house was quiet, and there was nothing else to take off any fire.

I got up on the arm of her chair, which I am not allowed to do, and nobody told me to get down. Her hands were folded in her lap, and her face was not frightened, and not surprised, and not anything. Between a kettle boiling and a kettle screaming is no time at all. Whoever did this to her did it in that, in her own parlor, in her own light: faster than any needle honestly goes.

Fifteen years. Every evening. She always came to the window.

I do not know how to be a cat that nobody is expecting.

I stayed on the arm of the chair a long time, and I will not tell you how long, because I do not know. The shutters stayed shut. At some point I found myself looking at the sewing table. Her red kerchief lay on it, folded, washed that morning, waiting to go back around her neck, where it had lived so long it was less a garment than a fact of her.

I put my head through it. Nobody tied it. It stays.

The thread at my chest, the tether every familiar wears, hers, cut like everything else in this house.

Not quite.

One strand left. Thin as spite. Running out from under the window, over the sill where her geraniums used to argue with the slugs, down into the fog, into the city.

I looked at it a long time. Outside, eight bells rang somewhere, and no ninth.

Then that is where we are going.
